\documentclass[../../main2.tex]{subfiles}

\begin{document}

\chapter{数学附录：把公式捏在手里}
\label{app:v2math}

\begin{motivation}
	你可能会问：正文不是说清楚了吗，为什么还要数学附录？\\
	\textbf{核心动机}：正文给直觉，附录给可检验的承诺。「贡献是条件概率差」这句话，要说清它依赖什么假设、放松后哪里失效、多因时怎么算。形式化不是炫耀，是让「大概如此」变得\textbf{可争议}。
\end{motivation}

\begin{logicmap}
	符号与生活映射（A.0）$\rightarrow$ 条件概率（A.1）$\rightarrow$ 链式分解（A.2）$\rightarrow$ 贡献分解（A.3）$\rightarrow$ 交互项与分账（A.4）$\rightarrow$ 归一化（A.5）$\rightarrow$ 边界与失效（A.6）
\end{logicmap}

\section{A.0 怎么用这份附录}

三条读法，先说好：

\begin{itemize}
	\item \textbf{按需跳读}：每一节标了对应正文小节，读懂哪节回哪节；
	\item \textbf{每个定义都带「为什么这样定义好」}：不只给定义，还告诉你它比朴素直觉多排除了什么歧义；
	\item \textbf{每个公式都有能算出数的算例}：检验标准只有一个——完全不懂统计的读者，能不能照例子自己算一道。
\end{itemize}

\textbf{符号先给生活映射}（这一张表是本附录的地基，忘词先回来看）：

\begin{center}
\begin{tabular}{lll}
\toprule
符号 & 读作 & 生活上对应什么 \\
\midrule
$P(X)$ & 事件 $X$ 的概率 & 「一百回里有几回」 \\
$P(D \mid A)$ & 条件概率 & 「$A$ 发生时，$D$ 一百回里有几回」 \\
$\neg A$ & 非 $A$ & 「没有 $A$ 的那个平行宇宙」 \\
$C_A$ & $A$ 的贡献 & 「拿掉 $A$，成功率掉几成」 \\
$I_{AB}$ & 交互项 & 「两人合谋多出来的（或抢功少掉的）那块」 \\
状态空间 & 所有可能处境的集合 & 「一个人所有可能摊上的事」 \\
$u_i(x_i)$ & 第 $i$ 把尺子 & 「体检表上的一项」 \\
$w_i$ & 权重 & 「你心里各项的相对分量」 \\
\bottomrule
\end{tabular}
\end{center}

\begin{questionbox}
	\textit{现在你可能想反驳：「正文不是说清楚了吗，形式化是不是在炫技？」}\\
	\textbf{不是。}正文说「排污占 25\%」，附录说清「这个 25\% 在哪四个假设下才成立」。正文给结论，附录给\textbf{可攻击面}——没有可攻击面的知识，不是知识，是安慰剂。
\end{questionbox}

\paragraph*{逻辑衔接}
	地基铺好。先立第一个正式对象：条件概率——「在……的条件下」这五个字的精确化。

\section{A.1 条件概率：定义与方向}
\label{app:condprob}

\begin{definitionbox}
	\textbf{定义}（条件概率）。对两个事件 $A, B$，且 $P(B) > 0$：
	\[
	P(A \mid B) \;=\; \frac{P(A \cap B)}{P(B)}.
	\]
	读作「在 $B$ 发生的条件下，$A$ 的概率」。
\end{definitionbox}

\textbf{为什么这样定义好？}朴素的「A 和 B 有关系」是感觉，这个式子把「有关系」钉成\textbf{可以排队的数}：两个不同的 $B$ 条件下，$A$ 的概率谁大谁小，一目了然。而且分母 $P(B)$ 明确了「你站在哪群人里看」——同一个事件，换人群概率就换，这件事不写清楚，后面全会乱。

\begin{mathbox}{算例：体检异常与真患病}
	1000 人体检，100 人真患病，其中 90 人指标异常；未患病的 900 人里，180 人也异常。\\
	\textbf{第 1 步}：$P(\text{异常}) = (90+180)/1000 = 0.27$。\\
	\textbf{第 2 步}：$P(\text{患病} \mid \text{异常}) = 90 / 270 = 0.333$。\\
	\textbf{第 3 步}：$P(\text{异常} \mid \text{患病}) = 90 / 100 = 0.90$。\\
	\textbf{对账}：0.90 和 0.333 差了近三倍——「异常时多大概率病」和「病了多大概率异常」是两个宇宙的问题。\end{mathbox}

\textbf{因果方向声明}：全书只用 $P(\text{果} \mid \text{因})$（从因推果）。诊断方向 $P(\text{因} \mid \text{果})$ 只在显式声明先验、经贝叶斯翻转时出场——两个方向混用，正是旧稿附录断裂的地方，本版从源头堵死。

\begin{mathbox}{贝叶斯翻转（备用工具，用时必须报先验）}
	\[
	P(A \mid B) = \frac{P(B \mid A)\, P(A)}{P(B)}.
	\]
	\textbf{算例}（接上题）：$P(\text{患病}) = 0.1$，$P(\text{异常} \mid \text{患病}) = 0.9$，$P(\text{异常}) = 0.27$，则
	$P(\text{患病} \mid \text{异常}) = (0.9 \times 0.1) / 0.27 = 0.333$——与上题对上账。\\
	\textbf{为什么这样定义好？}翻转公式强迫你报出先验 $P(A)$——「先验从哪来」永远摆在台面上，没法偷渡。
\end{mathbox}

\paragraph*{逻辑衔接}
	条件概率能算「一步」了。但预测常常是接力：一环扣一环。下一节把长链拆成可测量的接口。

\section{A.2 条件概率链的乘积分解}
\label{app:chain}

\mathlink{sec:pred-prob}{正文链接：第 6 章「两种预测」}

\begin{definitionbox}
	\textbf{链式假设}（马尔可夫性）：因果链 $C_0 \to C_1 \to \cdots \to C_n \to R$ 上，每一环只依赖上一环——「隔代不直接说话」。
\end{definitionbox}

\begin{mathbox}{链式分解与推导}
	在链式假设下：
	\[
	P(R \mid C_0) \;=\; \prod_{k=1}^{n} P(C_k \mid C_{k-1}) \cdot P(R \mid C_n).
	\]
	\textbf{推导}：联合事件的乘法法则给出
	$P(R \mid C_0) = P(C_1 \mid C_0) \, P(C_2 \mid C_1) \cdots P(R \mid C_n)$——每一步都是把「长链」剪成「相邻两环」；链式假设保证剪断处不留残余联系。\\
	\textbf{为什么这样定义好？}每个因子都是\textbf{可独立测量的接口}：供应链上每道工序测良率，政策链上每个环节测通过率——长问题变成一排短问题。
\end{mathbox}

\begin{mathbox}{算例：供应链三道工序}
	来料合格率 $0.90$，加工良率 $0.80$，出厂检出率 $0.95$。\\
	\textbf{第 1 步}：$P(\text{合格出厂} \mid \text{来料}) = 0.90 \times 0.80 \times 0.95 = 0.684$。\\
	\textbf{第 2 步}：短板在哪？$0.80$ 最低——加工环节是杠杆点。把它提到 $0.90$，整链变 $0.90 \times 0.90 \times 0.95 = 0.7695$，涨 8.6 个百分点；把最高的 $0.95$ 提到 $1.0$，只涨 3.5 个点。\\
	\textbf{第 3 步}：冗余怎么算？加工环节并联一台备用机（两台全坏才算坏，坏率 $0.2 \to 0.2^2 = 0.04$），整链变 $0.90 \times 0.96 \times 0.95 \approx 0.82$。\end{mathbox}

\textbf{边界}：链式假设是\textbf{假设}，不是定理。现实中「隔代说话」（跳环的因果）很常见——漏一条旁路，乘积就偏。用这个分解时必须自问：有没有绕过中间环节的暗线？

\paragraph*{逻辑衔接}
	单链会算了。但「拿掉 A 掉多少」这个全书的核心量，还没正式登场。下一节：贡献分解。

\section{A.3 贡献分解：条件概率差}
\label{app:contrib}

\mathlink{sec:butfor-formula}{正文链接：第 8 章「同一动作，写成公式」}

\begin{definitionbox}
	\textbf{定义}（贡献）。结果 $D$、候选原因 $A$：
	\[
	C_A \;=\; P(D \mid A) \;-\; P(D \mid \neg A).
	\]
	读作：「有 $A$ 的世界减没有 $A$ 的世界，差值就是 $A$ 的贡献」。
\end{definitionbox}

\textbf{为什么这样定义好？}三件事一次到位：① 把「我觉得他有责任」钉成可比较的量；② 天然含\textbf{平行宇宙对比}——正是无你检验的操作；③ 基准线 $P(D \mid \neg A)$ 把「本来就会发生」的部分扣掉，功劳不冒领。

\begin{mathbox}{算例：排污与鱼塘（正文第 8 章案例的完整版）}
	$P(\text{鱼苗死} \mid \text{排污}) = 0.30$，$P(\text{鱼苗死} \mid \text{不排污}) = 0.05$。\\
	\textbf{第 1 步}：$C_{\text{排污}} = 0.30 - 0.05 = 0.25$。\\
	\textbf{第 2 步}（读数）：排污把死亡率抬高 25 个百分点——这 25\% 记它的账。\\
	\textbf{第 3 步}（比较）：另一家工厂 $C = 0.12$——出事概率差一倍，责任可以排队了。\\
	\textbf{第 4 步}（反算）：若鉴定测得 $P(D \mid A) = 0.05 = P(D \mid \neg A)$，则 $C_A = 0$——排污期和对照期一样，账不成立，判不下去。
\end{mathbox}

\begin{mathbox}{算例：迟到与奖金（日常尺度）}
	$P(\text{丢奖} \mid \text{迟到}) = 0.8$，$P(\text{丢奖} \mid \text{不迟到}) = 0.1$，$C_{\text{迟到}} = 0.8 - 0.1 = 0.7$。\\
	\textbf{自查}：若改口径「丢的是一整年奖金」，$P$ 变，$C$ 就变——贡献\textbf{永远是针对某个口径的 D}，说 $C$ 之前先说清 $D$。
\end{mathbox}

\textbf{边界三条}：① $C_A$ 只对「可拿掉」的东西定义（基因、历史拿不掉，不进公式）；② $C_A$ 是概率差，不是道德分（台风也有 $C$）；③ 两个世界只有一个能观测——$C$ 永远是\textbf{估}出来的，第 7 章三把镊子是估它的工具箱。

\paragraph*{逻辑衔接}
	单因会算了。但两因同台时，加不起来的部分去哪了？下一节：交互项——全附录最硬的一块。

\section{A.4 多因一果：交互项与分账}
\label{app:interaction}

\mathlink{sec:multiple-interaction}{正文链接：第 9 章「交互项：一加一不等于二」}

\begin{definitionbox}
	\textbf{定义}（交互项）。两个原因 $A, B$ 同时作用于 $D$：
	\[
	I_{AB} = \big[P(D \mid A,B) - P(D \mid \neg A, \neg B)\big] - C_A - C_B,
	\]
	其中 $C_A = P(D \mid A,B) - P(D \mid \neg A, B)$，$C_B = P(D \mid A,B) - P(D \mid A, \neg B)$（边际定义，见算例）。
\end{definitionbox}

\textbf{为什么这样定义好？}它不许你和稀泥：$C_A + C_B + I_{AB}$ 精确等于总抬升——多不多少不少，合谋的那块\textbf{单列}。谁想把它偷偷塞给 A 或 B，账面立刻对不上。

\begin{mathbox}{算例：工地事故（正文第 9 章数字的完整对账）}
	$P(D \mid A,B) = 0.30$，$P(D \mid \neg A, \neg B) = 0.02$，$P(D \mid A,\neg B) = 0.12$，$P(D \mid \neg A, B) = 0.08$。\\
	\textbf{第 1 步}：总抬升 $= 0.30 - 0.02 = 0.28$。\\
	\textbf{第 2 步}：$C_A = 0.30 - 0.08 = 0.22$（有 B 时拿掉 A，看差多少）。\\
	\textbf{第 3 步}：$C_B = 0.30 - 0.12 = 0.18$。\\
	\textbf{第 4 步}：$I_{AB} = 0.28 - 0.22 - 0.18 = -0.12$——负的！防护网与堆载功能重叠，互相抢功。\\
	\textbf{第 5 步}（对账）：$0.22 + 0.18 + (-0.12) = 0.28$ $\checkmark$ 分文不差。
\end{mathbox}

\begin{mathbox}{算例：夏普利值分账（两方，可手算）}
	设事故总损失可归因的「价值」为 $v$：$v(\varnothing)=0$，$v(\{A\})=60$ 万，$v(\{B\})=40$ 万，$v(\{A,B\})=120$ 万。\\
	\textbf{第 1 步}（A 的夏普利值）：把上场顺序排一遍——A 先上：贡献 $v(\{A\})-v(\varnothing)=60$；A 后上：贡献 $v(\{A,B\})-v(\{B\})=80$。平均：$\varphi_A = \tfrac{60+80}{2} = 70$ 万。\\
	\textbf{第 2 步}（B 同理）：$\varphi_B = \tfrac{40+60}{2} = 50$ 万。\\
	\textbf{第 3 步}（对账）：$70 + 50 = 120$ 万 $\checkmark$ 交互的 20 万被两人均摊（各 10 万）。\\
	\textbf{第 4 步}（对照边际法）：$C_A = 120-40 = 80$，$C_B = 120-60 = 60$，合计 140 万 $> 120$ 万——边际法双记了交互项，所以\textbf{只能排序，不能当百分比}。
\end{mathbox}

\textbf{怎么选}：要分钱（必须加总）用夏普利；要抓主因（定整改优先级）用边际。两套都是约定，选了就认账——这是第 2 章不可通约性在因果空间的回声。

\paragraph*{逻辑衔接}
	分账工具齐了。但把贡献变成「百分比」之后有个陷阱——最后一节数学，也是最后一处诚实标注。

\section{A.5 归一化与量级}
\label{app:normalize}

\begin{mathbox}{算例：份额一样，风险差十倍}
	小事故：$C_A=0.04, C_B=0.06$；大灾难：$C_A=0.40, C_B=0.60$。\\
	\textbf{第 1 步}：份额都是 $40\% : 60\%$——两张份额表一模一样。\\
	\textbf{第 2 步}：总量差十倍（$0.10$ 对 $1.00$）。只看份额，你会对大灾难按下同样的重视程度。\\
	\textbf{规矩}：$s_A = C_A / \sum_i C_i$ 可以报，但必须与 $\sum_i C_i$ 成对出现。
\end{mathbox}

\textbf{为什么这样定义好（针对份额）？}归一化让不同规模的案例可比——代价是丢掉量级。定义的「好」从来不是免费的：知道它丢了什么，才知道什么时候用它。

\section{A.6 边界、假设与失效清单}

这份附录的全部承诺，站在四条假设上。哪条松了，哪里塌：

\begin{enumerate}
	\item \textbf{链式假设}（A.2）：隔代有暗线，乘积分解偏。塌法：整链概率高估。
	\item \textbf{可拿掉}（A.3）：$A$ 必须可干预。不可干预的东西（基因、历史）没有 $C_A$——别硬算。
	\item \textbf{二值清晰的事件}（A.1–A.4）：$A$、$D$ 要有明确的「发生/不发生」。模糊口径（「内卷」算不算发生了？）会让所有数失去意义。
	\item \textbf{平稳性}：估出的 $P$ 来自历史。环境变，$P$ 就变（正文第 15 章失败模式 2）。
\end{enumerate}

\textbf{我们能说什么}：给定口径、给定假设，$C_A$ 与 $I_{AB}$ 是可争议、可复算的量；不同案例可以排队比较。\\
\textbf{我们不能说什么}：$C$ 的「真值」；不可干预之物的贡献；预测本身在回音最大的地方（第 16 章）的精度。

\begin{warningbox}
	本附录与正文冲突时，以本附录的严格版本为准——直觉向严格让步。但若本附录与\textbf{现实}冲突，以现实为准——数学向事实让步。一本讲预测的书，自己先不诚实，就完了。
\end{warningbox}

\begin{reader_anticipation}
	\item 这些公式和统计课上的是一回事吗？→ 内核一样，用途不同：课本估「平均效应」，本书盯「拿掉谁」——反事实对比是主语。
	\item 想看这些公式怎么装回社会？→ 回正文第 11–15 章：四条通道逐个代入，第 15 章拼回整表。
\end{reader_anticipation}

\begin{thinkbox}
\begin{enumerate}
	\item 照 A.3 的步骤，自己算一道：某「早起」习惯与「全天精神」的关系——$P(\text{精神} \mid \text{早起}) = 0.6$，$P(\text{精神} \mid \text{不早起}) = 0.35$。$C$ 是多少？再说说哪个混杂因子（睡眠时长？）可能污染它。
	\item \textbf{反事实}：若 A.2 的三道工序里，两台「并联备用」的坏率\textbf{不独立}（比如停电会同时坏），第 3 步的 $0.2^2 = 0.04$ 还成立吗？整链会变多少？——这是「假设松了哪里塌」的现场练习。
\end{enumerate}
\end{thinkbox}

\end{document}
